% sec-05: outline item E.  Figures 13 (mermaid), 14 (plantuml), 15 (graphviz).
% Tables 14 and 15 are keyed by the same twelve IDs.
% FINAL stage improvements: Figure 14 gains a branch owner on each column title;
% Figure 15 gains a gate-type legend.
% Figure 13 no longer carries the author's \DiagramXShift adjustwidth wrapper,
% correction 11 of useredits.md.  The wrapper was compensating for a defect in
% capstyle.sty, where a caption or a frame ended its paragraph with glue that
% \par then deleted, driving every float 13.1pt right of centre.  That is fixed
% at source.  The wrapper was then measured at 0mm and still moved this figure
% 1.86pt left of every other one, because adjustwidth is a list environment and
% a list carries its own margin handling; removing it puts Figure 13 on 306.00pt
% with the other nineteen.  The carrier itself is kept, documented and set to
% 0mm at the foot of capstyle.sty, for the next figure that genuinely needs an
% optical nudge.
\section{Twelve Milestones a Program Officer Can Audit}

\subsection{What an Auditable Milestone Is}

A milestone is auditable only if it carries four things: an evidence artifact
that exists as a file, a cost that sums into its phase, a stop condition stated
before the work begins, and a named funder mechanism that pays for it. All
twelve rows below carry all four. The five Phase I rows sum to \$306,000 and the
seven Phase II rows sum to \$1,300,000, both exactly, and a milestone table that
does not sum to its award is the fastest way to lose a reader who has read
several hundred of them. Figure 13 places all twelve on the thirty-three month
calendar with the artifact date marked on each.

\begin{appfloat}[!tb]
\begin{appfig}[mermaid-type / gantt]
% 33 months at 0.415cm per month.  Row pitch 0.52cm against a 0.36cm bar
% height leaves 1.6mm of clear space between rows.  Row labels sit to the
% right of their own bar and never above it.
\foreach \m/\lab in {0/{month 0}, 6/{6}, 12/{12}, 18/{18}, 24/{24}, 30/{30}, 33/{33}}{
  \draw[pagraym,line width=0.3pt] (\m*0.415,0.36) -- (\m*0.415,-6.35);
  \node[font=\tiny\sffamily,text=protogray,anchor=south] at (\m*0.415,0.39) {\lab};}
\foreach \i/\y/\s/\w/\lab/\cost in {%
  1/0/1/1/{M1 site feasibility executed}/{\$24,000},
  2/-0.52/2/2/{M2 IRB submission accepted}/{\$31,000},
  3/-1.04/3/3/{M3 interlock bench verification}/{\$96,000},
  4/-1.56/4/3/{M4 VVUQ suite frozen and hashed}/{\$73,000},
  5/-2.08/6/3/{M5 IND amendment safe to proceed}/{\$82,000}}{
  \node[mmbar,minimum width={\w*0.415cm},anchor=west] at ({\s*0.415},\y) {};
  \node[mmdec,text width=2mm,minimum height=1mm,scale=0.30] at ({(\s+\w)*0.415},\y) {};
  \node[font=\tiny\sffamily,anchor=west] at ({(\s+\w)*0.415+0.16},\y) {\lab};
  \node[d2cell,minimum width=15mm,text width=12mm,minimum height=4.2mm,anchor=east]
       at (14.35,\y) {\cost};}
\draw[pagrayd,line width=0.5pt] (0,-2.47) -- (14.35,-2.47);
\node[font=\tiny\sffamily\bfseries,anchor=east] at (14.35,-2.35) {Phase I total \$306,000};
\foreach \i/\y/\s/\w/\lab/\cost in {%
  6/-2.86/10/3/{M6 first participant dosed}/{\$164,000},
  7/-3.38/12/4/{M7 first advised robotic Whipple}/{\$228,000},
  8/-3.90/15/5/{M8 dose level 1 cleared, n is 3}/{\$196,000},
  9/-4.42/18/6/{M9 audit replay to a third party}/{\$131,000},
  10/-4.94/22/6/{M10 dose level 2 cleared, n is 6}/{\$242,000},
  11/-5.46/26/5/{M11 interim PK, PD and ctDNA}/{\$187,000},
  12/-5.98/30/3/{M12 closeout and public deposit}/{\$152,000}}{
  \node[mmbark,minimum width={\w*0.415cm},anchor=west] at ({\s*0.415},\y) {};
  \node[mmdec,text width=2mm,minimum height=1mm,scale=0.30] at ({(\s+\w)*0.415},\y) {};
  \node[font=\tiny\sffamily,anchor=west] at ({(\s+\w)*0.415+0.16},\y) {\lab};
  \node[d2cell,fill=pablue2,minimum width=15mm,text width=12mm,minimum height=4.2mm,anchor=east]
       at (14.35,\y) {\cost};}
\draw[protoblue,line width=1pt,dashed] (3.735,0.55) -- (3.735,-6.35);
\node[font=\tiny\sffamily\bfseries,text=protoblue,anchor=south] at (3.735,0.58) {gate, month 9};
\node[font=\tiny\sffamily\bfseries,anchor=east] at (14.35,-6.32) {Phase II total \$1,300,000};
\node[mmlanetitle,anchor=east,align=right] at (-0.15,-1.04) {Phase I};
\node[mmlanetitle,anchor=east,align=right] at (-0.15,-4.42) {Phase II};
\node[font=\tiny\itshape,text=protogray,anchor=north west,text width=132mm]
  at (-1.35,-6.80) {The diamond at the right end of each bar is the date the
  evidence artifact exists, not the date the work stops. The one-month dead band
  at the gate is real: M5 ends at month 9 and M6 begins at month 10.};
\end{appfig}
\vspace{-0.65cm}
\figcaption{Figure 13. Thirty-three months, twelve milestones, and the artifact\\
date on each. Phase I holds five at 306,000 dollars; Phase II holds\\
seven at 1,300,000. The gate at month nine is the only hard boundary.}
\end{appfloat}

\subsection{The Twelve Rows}

The rows are given twice. Table 14 carries the milestone, its month and its
cost, and Table 15 keys the same twelve to what proves each one, what halts it,
and which mechanism pays for it.

\begin{apptable}
\begin{tabularx}{\textwidth}{>{\raggedright\arraybackslash}p{0.9cm} >{\raggedright\arraybackslash}p{1.5cm} >{\raggedright\arraybackslash}X >{\raggedright\arraybackslash}p{1.7cm} >{\raggedright\arraybackslash}p{1.6cm}}
\headrow \thead{ID} & \thead{Months} & \thead{Milestone} & \thead{Cost} & \thead{Phase}\\
\midrule
M1 & 1 to 2 & Site feasibility executed & \$24,000 & Phase I \\
M2 & 2 to 4 & IRB submission accepted & \$31,000 & Phase I \\
M3 & 3 to 6 & Interlock rig bench verification & \$96,000 & Phase I \\
M4 & 4 to 7 & VVUQ suite frozen and hashed & \$73,000 & Phase I \\
M5 & 6 to 9 & IND amendment safe to proceed & \$82,000 & Phase I \\
M6 & 10 to 13 & First participant dosed & \$164,000 & Phase II \\
M7 & 12 to 16 & First advised robotic Whipple completed & \$228,000 & Phase II \\
M8 & 15 to 20 & Dose level 1 cleared, three participants & \$196,000 & Phase II \\
M9 & 18 to 24 & Audit replay demonstrated to a third party & \$131,000 & Phase II \\
M10 & 22 to 28 & Dose level 2 cleared, six participants & \$242,000 & Phase II \\
M11 & 26 to 31 & Interim PK, PD and ctDNA readout & \$187,000 & Phase II \\
M12 & 30 to 33 & Closeout and public deposit & \$152,000 & Phase II \\
\end{tabularx}
\end{apptable}
\vspace{-0.65cm}
\tabcap{Table 14. The twelve milestones with their months and cost. M1\\
to M5 sum to \$306,000 and M6 to M12 sum to \$1,300,000; the same\\
twelve costs appear as work packages WP1 to WP12 in Figure 9.}

\begin{apptable}
\begin{tabularx}{\textwidth}{>{\raggedright\arraybackslash}p{0.9cm} >{\raggedright\arraybackslash}p{4.2cm} >{\raggedright\arraybackslash}X >{\raggedright\arraybackslash}p{2.1cm}}
\headrow \thead{ID} & \thead{Evidence artifact} & \thead{Stop condition} & \thead{Mechanism}\\
\midrule
M1 & Countersigned feasibility memorandum & No site by month 4; Phase I withdrawn & SBIR Phase I \\
M2 & IRB acknowledgment and protocol v1.0 & Two rejections; protocol re-scoped & SBIR Phase I \\
M3 & 200-run latency report with p95 table & p95 above 250 ms; no first-in-human & SBIR Phase I \\
M4 & VVUQ report and SHA-256 manifest & Any credibility factor below its gate & SBIR Phase I \\
M5 & FDA acknowledgment, 30-day clock closed & Clinical hold; Phase II not requested & SBIR Phase I \\
M6 & Consent, dispensing record, day-1 CRF & No enrollment by month 15; descope to three & SBIR Phase II \\
M7 & Operative note, advisory log, replay bundle & Any advisory override event; advisory read-only & SBIR Phase II \\
M8 & DSMB minutes and safety table & Two of three DLTs; de-escalate & SBIR Phase II \\
M9 & Replay record and hash-match certificate & Any non-reproducible record; hold enrollment & SBIR Phase II \\
M10 & Cumulative safety table, ISGPS grades & Grade B or C fistula above 30 percent & SBIR Phase II \\
M11 & PK report and ctDNA clearance table & No target engagement at dose level 2 & SBIR Phase II \\
M12 & Zenodo deposit, report, repository tag & None; this is the terminal artifact & SBIR Phase II \\
\end{tabularx}
\end{apptable}
\vspace{-0.65cm}
\tabcap{Table 15. The same twelve rows keyed by ID, carrying what proves\\
each milestone, what halts it, and which mechanism pays. Every stop\\
condition is a number or a signed instrument; none is a judgement call.}

M6 additionally requires that the trial be registered before the first
participant is enrolled, and the registration record is part of that milestone's
artifact set \cite{ctepiit}.

Everything in the \$2,104,000 delta of \S3 lies beyond M12 and carries no
milestone row here, deliberately. A milestone schedule that promised work no
award pays for would not be auditable.

\subsection{How a Milestone Is Closed}

Figure 14 runs the three evidence branches of a single milestone concurrently
and joins them before a program officer can audit anything.

\begin{appfloat}[!tb]
\begin{appfig}[plantuml-type / activity with fork and join]
% Top down, because a fork and join is genuinely hierarchical in time.  Branch
% pitch 5.35cm horizontal, 1.00cm vertical.  The one return edge is routed at
% x = -0.30, outside the leftmost branch column.
\node[umlinit] (p0) at (6.90,0.35) {};
\node[umlbox,text width=34mm] (p1) at (6.90,-0.40) {milestone work executes};
\node[umlbar,minimum width=88mm] (fk) at (6.90,-1.30) {};
\foreach \i/\y/\lab in {1/-2.10/{company produces artifact}, 2/-3.10/{hash with SHA-256},
  3/-4.10/{deposit to repository}, 4/-5.10/{tag the release}}{
  \node[umlbox,text width=27mm] (A\i) at (1.55,\y) {\lab};}
\foreach \i/\y/\lab in {1/-2.10/{site produces source record}, 2/-3.10/{investigator signs it},
  3/-4.10/{retain in the site file}}{
  \node[umlbox,text width=27mm] (B\i) at (6.90,\y) {\lab};}
\foreach \i/\y/\lab in {1/-2.10/{independent monitor visits}, 2/-3.10/{verify source to artifact},
  3/-4.10/{write monitoring report}}{
  \node[umlbox,text width=27mm] (C\i) at (12.25,\y) {\lab};}
\node[d2title,anchor=south] at (1.55,-1.72) {A, the company};
\node[d2title,anchor=south] at (6.90,-1.72) {B, the site};
\node[d2title,anchor=south] at (12.25,-1.72) {C, an independent monitor};
\node[umlbar,minimum width=88mm] (jn) at (6.90,-5.90) {};
\node[umlkey,text width=34mm] (p2) at (6.90,-6.70) {artifact bundle complete};
\node[mmdec,text width=20mm] (d1) at (6.90,-7.95) {audit requested?};
\node[mmdec,text width=20mm] (d2) at (3.10,-9.35) {hash matches?};
\node[umlstateon,text width=26mm] (ok) at (10.90,-9.35) {milestone accepted};
\node[umlstategray,text width=26mm] (re) at (0.55,-10.65) {milestone reopened};
\node[umlfinal] (p9) at (13.45,-10.65) {};
\draw[umlarrow] (p0) -- (p1);
\draw[umlarrow] (p1) -- (fk);
\foreach \x in {1.55,6.90,12.25}{\draw[umlarrow] (\x,-1.36) -- (\x,-1.90);}
\draw[umlarrow] (A1) -- (A2); \draw[umlarrow] (A2) -- (A3); \draw[umlarrow] (A3) -- (A4);
\draw[umlarrow] (B1) -- (B2); \draw[umlarrow] (B2) -- (B3);
\draw[umlarrow] (C1) -- (C2); \draw[umlarrow] (C2) -- (C3);
\draw[umlarrow] (1.55,-5.42) -- (1.55,-5.84);
\draw[umlarrow] (6.90,-4.42) -- (6.90,-5.84);
\draw[umlarrow] (12.25,-4.42) -- (12.25,-5.84);
\draw[umlarrow] (jn) -- (p2);
\draw[umlarrow] (p2) -- (d1);
\draw[umlarrow] (d1.west) -- node[umlguard,pos=0.5] {[yes]} (d2.north east);
\draw[umlarrow] (d1.east) -- node[umlguard,pos=0.5] {[no, next report]} (ok.north west);
\draw[umlarrow] (d2.east) -- node[umlguard,pos=0.45] {[match]} (ok.west);
\draw[umlarrow] (d2.west) -- node[umlguard,pos=0.5] {[no match]} (re.north);
\draw[umldash] (re.west) to[bend left=36]
  node[umlguard,pos=0.5] {[rework]} (p1.west);
\draw[umlarrow] (ok) -- (p9);
\node[font=\tiny\itshape,text=protogray,anchor=north west,text width=132mm]
  at (-0.30,-11.55) {All three branches must join before an audit can begin,
  because an audit compares them: the hash proves the artifact has not moved,
  the source record proves it describes something that happened, and the
  monitoring report proves a third party looked at both.};
\end{appfig}
\vspace{-0.65cm}
\figcaption{Figure 14. One milestone, three concurrent evidence branches, and the\\
join that has to close before a program officer can audit anything.\\
The replay is the only step that can send a milestone backwards.}
\end{appfloat}

Any one branch alone is a claim rather than evidence. A hash with no source
record proves that a file has not changed and says nothing about whether it
describes an event. A source record with no hash proves an event occurred and
says nothing about whether the artifact still matches it. The monitoring report
is what joins the two, and it is written by somebody the company does not employ
\cite{ecfr2026part312records}.

\subsection{How the Programme Stops}

Figure 15 puts one top event over five halt points and ten basic events, so a
reader can see what has to fail, and in what combination, to stop the work.

\begin{appfloat}[!tb]
\begin{appfig}[graphviz-type / fault tree]
% Strictly layered across five ranks.  No edge skips more than one rank.  The
% five halt-to-gate edges are fanned across a 1.2cm span on the top gate's
% underside so their arrowheads are 3mm apart.
\node[gvboxk,text width=46mm] (top) at (7.10,0) {Programme stops before M12 without a deposited archive};
\umlgateor{7.10}{-1.05}
\draw[gvedge] (7.10,-0.79) -- (top.south);
\foreach \i/\x/\lab in {1/0.75/{H1 gate halt, month 9}, 2/4.00/{H2 safety halt, month 20},
  3/7.10/{H3 capital halt, month 22}, 4/10.20/{H4 site halt, month 6},
  5/13.45/{H5 ownership halt, any month}}{
  \node[gvboxs,text width=22mm] (h\i) at (\x,-2.45) {\lab};}
\foreach \i/\dx in {1/-0.60, 2/-0.30, 3/0, 4/0.30, 5/0.60}{
  \draw[gvedge] (h\i.north) -- ({7.10+\dx},-1.35);}
\umlgateor{0.75}{-3.55}
\umlgateand{4.00}{-3.55}
\umlgateand{7.10}{-3.55}
\umlgateand{10.20}{-3.55}
\umlgateor{13.45}{-3.55}
\foreach \x in {0.75,4.00,7.10,10.20,13.45}{\draw[gvedge] (\x,-3.29) -- (\x,-2.72);}
\foreach \x/\a/\b in {0.75/{Bench p95 above 250 ms}/{IND clinical hold issued},
  4.00/{2 of 3 DLTs at dose level 1}/{De-escalation exhausts the range},
  7.10/{Phase II award not made}/{Seed round unsigned by month 22},
  10.20/{CTA not executed by month 6}/{No alternate San Diego site},
  13.45/{Founder falls below 50 percent}/{Investigator acquires 54.2 interest}}{
  \node[gvboxg,text width=19mm] at ({\x-0.95},-5.05) {\a};
  \node[gvboxg,text width=19mm] at ({\x+0.95},-5.05) {\b};
  \draw[gvedge] ({\x-0.95},-4.60) -- ({\x-0.28},-3.82);
  \draw[gvedge] ({\x+0.95},-4.60) -- ({\x+0.28},-3.82);}
\draw[pagrayd,line width=0.5pt] (-0.40,-6.05) -- (14.60,-6.05);
\umlgateand{1.60}{-7.75}
\node[font=\tiny\sffamily,text=protoblack,anchor=west] at (2.15,-7.90)
  {both basic events required};
\umlgateor{6.60}{-7.75}
\node[font=\tiny\sffamily,text=protoblack,anchor=west] at (7.15,-7.90)
  {either basic event suffices, a single point of failure};
\node[font=\tiny\sffamily\bfseries,text=protoblue,anchor=west] at (-0.40,-6.30)
  {Already deposited at each halt:};
\foreach \x/\lab in {0.75/{M1 to M4}, 4.00/{M1 to M8}, 7.10/{M1 to M9},
  10.20/{M1 partial}, 13.45/{all closed}}{
  \node[gvcellg,minimum width=24mm,text width=21mm,minimum height=5mm] at (\x,-6.85) {\lab};}
\node[font=\tiny\itshape,text=protogray,anchor=north west,text width=132mm]
  at (-0.40,-8.55) {The top event is a stop without a deposited archive, not a
  stop. Three of the five branches are AND gates. The two OR branches, H1 and
  H5, are the only single points of failure, and neither is technical: one is
  regulatory and instrumental, the other is a fact about who owns the shares.};
\end{appfig}
\vspace{-0.65cm}
\figcaption{Figure 15. One top event, five halt points, and ten basic events\\
below. Three are AND gates, so no single failure stops the programme.\\
The two OR branches are the month-nine gate and ownership test.}
\end{appfloat}

The two OR branches are the honest part of Figure 15. H1 fails on either a
bench latency above 250 ms at the 95th percentile or an IND clinical hold; H5
fails on either the founder falling below 50 percent or an investigator
acquiring a \S54.2 interest. Neither is technical in the sense that more
engineering fixes it, and both are knowable in advance, which is the only reason
they are acceptable to state on a milestone schedule.

The top event is deliberately not ``the programme stops''. It is a stop
\emph{without a deposited archive}, and the band beneath the tree is why the
distinction matters: at every halt point the archive is already deposited for
every milestone that has closed. Reaching the top event would require a halt
plus a separate failure to deposit, and depositing is itself a Phase I
deliverable that does not depend on the halt's cause. Figure 20 carries the
custody arrangement that makes that true.
